\documentclass[12pt]{article}

\usepackage[a4paper, top=1.5cm, bottom=1.5cm, left=1mm, right=2cm]{geometry}
\usepackage[english, russian]{babel} % для русского
\usepackage[T2A]{fontenc} % также для кириллицы
\usepackage{amsmath,amssymb} % для математических символов
\usepackage{graphicx} % Required for inserting images
\usepackage{multicol} % для нескольких колонок 
\usepackage{microtype} % для пробелов поменьше

\usepackage{tikz} % для самодельных рисунков


 
\begin{document}
    \pagestyle{empty}
    \setlength{\columnsep}{1pt} % расстояние между колонками

    \hspace{7cm} % контроль отступа слева
    {\centering 26 \qquad\qquad \textbf{КНИГА I ПРЕДЛ. II. ЗАДАЧА}}

    \vspace{8mm}
    \begin{multicols}{2}
    
    \begin{tikzpicture}
        \draw[color=red!100, ultra thick] circle (80pt);
        \draw[color=blue!100, ultra thick] (0.46,-0.98) circle (48pt);
        \draw[color=orange!100, ultra thick] (0.46, -0.98) -- (1.2, -2.47);
        \draw[densely dashed, color=red!100, ultra thick] (0.46, -0.98) -- (-0.04, 0.01);
        \draw[densely dashed, color=black!100, ultra thick] (0.46, -0.98) -- (1.1,-0.09);
        \draw[color=red!100, ultra thick] (-0.05, 0.02) -- (1.1,-0.09);
        \draw[color=black!100, ultra thick] (0.46, -0.98) -- (-1.2,-0.98);
        \draw[color=blue!100, ultra thick] (1.1,-0.09) -- (2.776, -0.27);
        \node[right, font=\tiny] at (2.776, -0.27) {F};
        \node[left, font=\tiny] at (-1.2, -0.98) {C};
        \node[below, xshift=3pt, font=\tiny] at (1.2, -2.47) {E};
        \node[below, xshift=-3pt, font=\tiny] at (0.46, -0.98) {B};
        \node[left, yshift=4pt, xshift=4pt, font=\tiny] at (-0.05, 0.01) {D};
        \node[above, xshift=3pt, yshift =-3pt, font=\tiny] at (1.1, -0.09) {A};
    \end{tikzpicture}

    \columnbreak % переходим на правую часть листа
    
    \noindent
    \begin{minipage}[t]{0.1\textwidth}
    \vspace{0pt}
    \includegraphics[width=\linewidth]{image.png}
    \end{minipage}
    % raisebox{0.5ex} - выравнивает область tikz рисунка, 0.5ex - типографская единица измерения, половина высоты "x"
    \noindent 
    \begin{minipage}[t]{0.7\textwidth}
    \vspace{0pt}
    \textit{т данной точки} \begin{tikzpicture} \raisebox{0.5ex}{\draw[color=red!100, ultra thick] (-0.6,0) -- (0,0);
    \draw[color=blue!100, ultra thick] (0,0) -- (0.6,0);}
    \node[above, font=\tiny] at (0,0) {A};
    \end{tikzpicture} \textit{отложить прямую,\\
    равную данной прямой} \begin{tikzpicture} \raisebox{0.5ex}{\draw[color=black!100, ultra thick] (0,0) -- (1.2,0);}
    \node[above, font=\tiny] at (0,0) {B}; \node[above, font=\tiny] at (1.2, 0) {C};
    \end{tikzpicture} .\\
    \end{minipage}

    \vspace{1cm}

    {\centering \noindent Проведем 
    \begin{tikzpicture} 
    \raisebox{0.5ex}{\draw[densely dashed, color=black!100, ultra thick](0,0) -- (1.2,0);}
    \node[above, font=\tiny] at (0,0) {A};
    \node[above, font=\tiny] at (1.2,0) {B};
    \end{tikzpicture} 
    (пост. I),\\
    построим \begin{tikzpicture} 
    \raisebox{-3.4ex}{\draw[densely dashed, color=red!100, ultra thick] (0,0) -- (0.52, -1);
    \draw[densely dashed, color=black!100, ultra thick] (1.2, -0.08) -- (0.52, -1); 
    \draw[color=red!100, ultra thick] (0,0) -- (1.2,-0.08);
    \node[above, xshift=-3, yshift=-3pt, font=\tiny] at (0,0) {D};
    \node[above, xshift=4pt, yshift=-4pt, font=\tiny] at (1.2, -0.08) {A};
    \node[below, yshift=2pt, font=\tiny] at (0.52, -1) {B};}
    \end{tikzpicture}
    (пр. I.i),\\
    \vspace{5mm}
    продлим \begin{tikzpicture} 
    \raisebox{0.5ex}{\draw[densely dashed, color=red!100, ultra thick] (0,0) -- (1.2,0); } 
    \node[above, font=\tiny] at (0,0) {B};
    \node[above, font=\tiny] at (1.2,0) {D};
    \end{tikzpicture} (пост. II),\\}
    
    \noindent опишем \begin{tikzpicture}
    \raisebox{-2.6ex}{\draw[color=blue!100, ultra thick] (0,0) circle (16pt);
    \draw[color=black!100, ultra thick] (0,0) -- (-0.54,0);
    \node[right, xshift=-3pt, font=\tiny] at (0,0) {B};
    \node[left, xshift=1pt, font=\tiny] at (-0.54, 0) {C};}
    \end{tikzpicture} (пост. III), и \begin{tikzpicture}
    \raisebox{-4.5ex}{\draw[color=red!100, ultra thick] (0,0) circle (26pt);
    \draw[densely dashed, color=red!100, ultra thick] (0,0) -- (0.15,-0.32);
    \draw[color=yellow!100, ultra thick] (0.15,-0.32) -- (0.38, -0.8);
    \node[above, xshift=-2pt, yshift=-3pt, font=\tiny] at (0,0) {D};
    \node[below, xshift=3pt, font=\tiny] at (0.38,-0.8) {E};
    } \end{tikzpicture} (пост. III);\\

    \vspace{4mm}
    
    {\centering продлим \begin{tikzpicture} 
    \raisebox{0.5ex}{\draw[color=red!100, ultra thick] (0,0) -- (1.2,0);}
    \node[above, font=\tiny] at (0,0) {D};
    \node[above, font=\tiny] at (1.2,0) {A};
    \end{tikzpicture} (пост. II),\\
    тогда искомая прямая это \begin{tikzpicture} \raisebox{0.5ex}{\draw[color=blue!100, ultra thick] (0,0) -- (1.2,0);} 
    \node[above, font=\tiny] at (0,0) {A};
    \node[above, font=\tiny] at (1.2,0) {F};
    \end{tikzpicture} .\\}

    \vspace{6mm} % небольшой отступ как в исходной странице

    {\centering \noindent Поскольку \begin{tikzpicture}
    \raisebox{0.5ex}{\draw[color=yellow!100, ultra thick] (0,0) -- (0.8,0);
    \draw[densely dashed, color=red!100, ultra thick] (0.8,0) -- (1.2,0);}
    \node[above, font=\tiny] at (0,0) {E};
    \node[above, font=\tiny] at (1.2,0) {D};
    \end{tikzpicture} $=$ \begin{tikzpicture}
    \raisebox{0.5ex}{\draw[color=red!100, ultra thick] (0,0) -- (0.4,0);
    \draw[color=blue!100, ultra thick] (0.4, 0) -- (1.2, 0);}
    \node[above, font=\tiny] at (0,0) {D};
    \node[above, font=\tiny] at (1.2,0) {F};
    \end{tikzpicture} (опр. \textit{15}),\\
    и \begin{tikzpicture}
    \raisebox{0.5ex}{\draw[densely dashed, color=red!100, ultra thick] (0,0) -- (1.2, 0);}
    \node[above, font=\tiny] at (0,0) {B};
    \node[above, font=\tiny] at (1.2,0) {D};
    \end{tikzpicture} $=$ \begin{tikzpicture}
    \raisebox{0.5ex}{\draw[color=red!100, ultra thick] (0,0) -- (1.2, 0);}
    \node[above, font=\tiny] at (0,0) {D};
    \node[above, font=\tiny] at (1.2,0) {A};
    \end{tikzpicture} (постр.),\\
    $\therefore$ \begin{tikzpicture}
    \raisebox{0.5ex}{\draw[color=yellow!100, ultra thick] (0,0) -- (1.2, 0);}
    \node[above, font=\tiny] at (0,0) {B};
    \node[above, font=\tiny] at (1.2,0) {E};
    \end{tikzpicture} $=$ \begin{tikzpicture}
    \raisebox{0.5ex}{\draw[color=blue!100, ultra thick] (0,0) -- (1.2, 0); }  
    \node[above, font=\tiny] at (0,0) {A};
    \node[above, font=\tiny] at (1.2,0) {F};
    \end{tikzpicture} (акс. III),\\
    но (опр. \textit{15}) \begin{tikzpicture}
    \raisebox{0.5ex}{\draw[color=black!100, ultra thick] (0,0) -- (1.2, 0);}
    \node[above, font=\tiny] at (0,0) {B};
    \node[above, font=\tiny] at (1.2,0) {C};
    \end{tikzpicture} $=$ \begin{tikzpicture}
    \raisebox{0.5ex}{\draw[color=yellow!100, ultra thick] (0,0) -- (1.2, 0);} 
    \node[above, font=\tiny] at (0,0) {B};
    \node[above, font=\tiny] at (1.2,0) {E};
    \end{tikzpicture} $=$ \begin{tikzpicture}
    \raisebox{0.5ex}{\draw[color=blue!100, ultra thick] (0,0) -- (1.2, 0);}
    \node[above, font=\tiny] at (0,0) {A};
    \node[above, font=\tiny] at (1.2,0) {F};
    \end{tikzpicture} .\\}

    \vspace{5mm} % ещё отступ

    $\therefore$ \begin{tikzpicture}
    \raisebox{0.5ex}{\draw[color=blue!100, ultra thick] (0,0) -- (1.2, 0);}
    \node[above, font=\tiny] at (0,0) {A};
    \node[above, font=\tiny] at (1.2,0) {F};
    \end{tikzpicture} проведенная из данной точки\\    
    (\begin{tikzpicture}
    \raisebox{0.5ex}{\draw[color=red!100, ultra thick] (0,0) -- (0.6, 0); 
    \draw[color=blue!100, ultra thick] (0.6,0) -- (1.2, 0);}
    \node[above, font=\tiny] at (0.6,0) {A};
    \end{tikzpicture}) равна данной прямой \begin{tikzpicture}
    \raisebox{0.5ex}{\draw[color=black!100, ultra thick] (0,0) -- (1.2, 0); }
    \node[above, font=\tiny] at (0,0) {B};
    \node[above, font=\tiny] at (1.2,0) {C};
    \end{tikzpicture} (акс. I).\\

    \hfill ч.т.д.
    
    \end{multicols}
    
\end{document}